% ============================================================
%  chapters/app-c-spherepop-calculus.tex
%  Appendix C: Spherepop Calculus as Nested Constraint Evaluation
% ============================================================

\section*{C.1 Introduction}

The Spherepop formalism arose from a simple geometric observation. Ordinary
symbolic notation represents hierarchy using brackets: $(1-(3\times 2^2))$.
The brackets indicate evaluation order but possess no geometric interpretation.
Spherepop replaces brackets by nested regions. Expressions become topological
objects. Evaluation becomes the progressive collapse of nested constraint domains.

% ── C.2–C.3 Expression Trees and Spherepop ────────────────────
\section*{C.2 Expression Trees}

\begin{definition}{Expression Tree}{exp-tree}
  Let $\Sigma$ be a set of symbols (constants, variables, operators). An
  \emph{expression tree} is a rooted tree $T = (V,E)$ whose leaves contain
  constants or variables and whose internal nodes contain operators.
\end{definition}

\begin{figure}[h]
  \centering
  \begin{tikzpicture}[scale=0.72]
    \tikzset{
  bubble/.style={circle, draw, thick, minimum size=7mm, inner sep=0},
  poparrow/.style={-{Stealth[length=7pt]}, RSVPaccent, very thick},
  steplab/.style={RSVPdeep, font=\small\bfseries\sffamily},
  nodelab/.style={RSVPdeep, font=\footnotesize},
  collapselab/.style={RSVPaccent!80, font=\footnotesize\itshape},
}
% ============================================================
%  diagrams/spherepop-eval-3d.tex
%  Spherepop evaluation sequence as nested 3D bubbles
%  Shows: 1 + 3×2² evaluated step by step
% ============================================================
  bubble/.style={circle, draw, thick, minimum size=7mm, inner sep=0},
  poparrow/.style={-{Stealth[length=7pt]}, RSVPaccent, very thick},
  steplab/.style={RSVPdeep, font=\small\bfseries\sffamily},
  nodelab/.style={RSVPdeep, font=\footnotesize},
  collapselab/.style={RSVPaccent!80, font=\footnotesize\itshape},
% ── STEP 0: Full nested expression 1 + 3×(2²) ──────────────
\begin{scope}[xshift=0cm, yshift=0cm]
  \node[steplab] at (0, 2.8) {Step 0: Initial expression};

  % Outermost bubble (the + operation)
  \draw[RSVPdeep, thick] (0,0) circle (2.4cm);
  \node[nodelab] at (-1.3, 1.5) {$+$};

  % Left atom: 1
  \draw[RSVPmid, thick] (-1.2, 0) circle (0.55cm);
  \node[nodelab] at (-1.2, 0) {$1$};

  % Middle bubble: × operation
  \draw[RSVPmid, thick] (0.8, 0) circle (1.4cm);
  \node[nodelab] at (-0.1, 0.8) {$\times$};

  % Left of ×: atom 3
  \draw[RSVPpurple, thick] (0.1, 0) circle (0.45cm);
  \node[nodelab] at (0.1, 0) {$3$};

  % Right of ×: inner bubble for 2²
  \draw[RSVPaccent, thick] (1.5, 0) circle (0.7cm);
  \node[nodelab] at (1.3, 0.35) {\footnotesize$\hat{\phantom{.}}$};
  % atom 2
  \draw[RSVPaccent!70, thick] (1.25, -0.1) circle (0.28cm);
  \node[font=\tiny] at (1.25,-0.1) {$2$};
  % atom 2 (exponent)
  \draw[RSVPaccent!70, thick] (1.7, 0.25) circle (0.22cm);
  \node[font=\tiny] at (1.7,0.25) {$2$};

  % Pop target indicator
  \draw[RSVPaccent, dashed, thick, opacity=0.6]
    (1.5, 0) circle (0.78cm);
  \node[collapselab] at (1.5, -1.05) {pop here};
\end{scope}

% ── Arrow 1 ─────────────────────────────────────────────────
\draw[poparrow] (2.8, 0) -- (4.2, 0)
  node[midway, above, collapselab] {$\bullet_1$};

% ── STEP 1: 2² evaluated → 4 ─────────────────────────────────
\begin{scope}[xshift=5.2cm, yshift=0cm]
  \node[steplab] at (0, 2.8) {Step 1: $2^2 \to 4$};

  \draw[RSVPdeep, thick] (0,0) circle (2.4cm);
  \node[nodelab] at (-1.3, 1.5) {$+$};

  \draw[RSVPmid, thick] (-1.2, 0) circle (0.55cm);
  \node[nodelab] at (-1.2, 0) {$1$};

  \draw[RSVPmid, thick] (0.8, 0) circle (1.4cm);
  \node[nodelab] at (-0.1, 0.8) {$\times$};

  \draw[RSVPpurple, thick] (0.1, 0) circle (0.45cm);
  \node[nodelab] at (0.1, 0) {$3$};

  % 2² has been evaluated: now a single node = 4
  \draw[RSVPaccent, very thick, fill=RSVPaccent!15] (1.5, 0) circle (0.5cm);
  \node[nodelab, font=\small\bfseries] at (1.5, 0) {$4$};

  % Pop next target: ×
  \draw[RSVPaccent, dashed, thick, opacity=0.6]
    (0.8, 0) circle (1.48cm);
  \node[collapselab] at (0.8, -1.55) {pop here};
\end{scope}

% ── Arrow 2 ─────────────────────────────────────────────────
\draw[poparrow] (8.0, 0) -- (9.4, 0)
  node[midway, above, collapselab] {$\bullet_2$};

% ── STEP 2: 3×4 → 12 ─────────────────────────────────────────
\begin{scope}[xshift=10.4cm, yshift=0cm]
  \node[steplab] at (0, 2.8) {Step 2: $3{\times}4 \to 12$};

  \draw[RSVPdeep, thick] (0,0) circle (2.4cm);
  \node[nodelab] at (-1.3, 1.5) {$+$};

  \draw[RSVPmid, thick] (-1.2, 0) circle (0.55cm);
  \node[nodelab] at (-1.2, 0) {$1$};

  % 3×4 evaluated → 12
  \draw[RSVPmid, very thick, fill=RSVPmid!15] (0.8, 0) circle (0.6cm);
  \node[nodelab, font=\small\bfseries] at (0.8, 0) {$12$};

  % Pop next: outer +
  \draw[RSVPaccent, dashed, thick, opacity=0.6]
    (0,0) circle (2.48cm);
  \node[collapselab] at (0, -2.55) {pop here};
\end{scope}

% ── Arrow 3 (below, wrapping) ────────────────────────────────
\draw[poparrow] (10.4, -3.1) -- (10.4, -4.0) -- (5.4, -4.0) -- (5.4, -3.2);
\node[collapselab] at (7.9, -4.3) {$\bullet_3$};

% ── STEP 3: 1+12 → 13 (final) ────────────────────────────────
\begin{scope}[xshift=5.4cm, yshift=-6.4cm]
  \node[steplab] at (0, 2.4) {Step 3: $1{+}12 \to 13$ \quad (collapsed)};

  % Final result: single atom
  \draw[RSVPdeep, very thick, fill=RSVPdeep!10] (0,0) circle (0.85cm);
  \node[RSVPdeep, font=\Large\bfseries] at (0,0) {$13$};

  \node[collapselab, font=\small] at (0, -1.2)
    {evaluation complete — depth 0};
\end{scope}

% ── Depth legend ─────────────────────────────────────────────
\begin{scope}[xshift=0cm, yshift=-5.5cm]
  \node[draw=RSVPdeep!30, fill=RSVPgray, rounded corners=3pt,
        inner sep=6pt, font=\scriptsize, anchor=north west] at (0,0) {
    \begin{tabular}{ll}
      \color{RSVPdeep}\rule{10pt}{1.5pt} & Outer sphere (depth 0–1)\\
      \color{RSVPmid}\rule{10pt}{1.5pt}  & Mid sphere (depth 1–2)\\
      \color{RSVPaccent}\rule{10pt}{1.5pt}& Inner sphere (depth 2–3)\\
      $\bullet_n$ & Pop at step $n$
    \end{tabular}
  };
\end{scope}

  \end{tikzpicture}
  \caption{Spherepop evaluation of $1 + 3 \times 2^2$. Each step collapses
    an innermost sphere (nested region). The complexity functional $C(E)$
    decreases monotonically: $3 \to 2 \to 1 \to 0$.}
  \label{fig:spherepop-eval}
\end{figure}

\begin{proposition}{Embedding Existence}{embedding-exists}
  Every finite expression tree admits a Spherepop embedding into nested
  closed regions.
\end{proposition}

\begin{proof}
  Proceed recursively. Embed the root as a region; embed each child as a
  disjoint interior region; repeat until leaves. Finiteness ensures termination.
\end{proof}

% ── C.5–C.6 Complexity and Pop ────────────────────────────────
\section*{C.5 The Complexity Functional}

\begin{definition}{Complexity Functional}{complexity}
  Let $E$ be an expression. Define $C(E)$ as the number of unresolved
  operator nodes.
\end{definition}

\begin{proposition}{Evaluation Decreases Complexity}{eval-decreases}
  Each evaluation step satisfies $C(E_{t+1}) = C(E_t) - 1$.
\end{proposition}

\begin{definition}{Pop Operator}{pop-def}
  The \emph{pop operator} $P$ maps an expression $E$ to a reduced
  expression $P(E)$ by eliminating one evaluation domain:
  \[
    P(2^2) = 4, \quad P(3 \times 4) = 12, \quad P(1 - 12) = -11.
  \]
\end{definition}

% ── C.9–C.10 Termination and Confluence ──────────────────────
\section*{C.9 Termination}

\begin{theorem}{Spherepop Termination}{sp-termination}
  Every finite arithmetic Spherepop expression terminates.
\end{theorem}

\begin{proof}
  Each pop decreases $C(E)$ by one. Since $C(E) \geq 0$ and is
  integer-valued, only finitely many decreases are possible.
\end{proof}

\section*{C.10 Confluence}

\begin{theorem}{Spherepop Confluence}{sp-confluence}
  Spherepop arithmetic is confluent whenever operator precedence is respected.
\end{theorem}

\begin{proof}
  Arithmetic evaluation defines a unique value. Every legal reduction sequence
  preserves semantic equivalence. Hence all legal paths terminate at the same value.
\end{proof}

% ── C.12–C.14 Constraint Descent and Renormalization ─────────
\section*{C.12 Spherepop as Constraint Descent}

Define $\kappa(E) = C(E)$ — the number of unresolved operations. Every pop
satisfies $\kappa(P(E)) < \kappa(E)$. Thus Spherepop evaluation is precisely
a constraint descent process (Appendix~A, Definition~\ref{def:constraint-descent}).

\begin{remark}
  This is the first major bridge between Spherepop and RSVP. Both are
  relaxation systems. Both reduce constraint violation. Both seek admissible
  configurations.
\end{remark}

\section*{C.14 Evaluation as Renormalization}

\begin{definition}{Renormalizing Pop}{renorm-pop}
  A \emph{renormalizing pop} is $R : M_i \to M_{i-1}$, replacing detailed
  internal structure by an effective value. Example: $(3 \times 2^2)$
  becomes $12$. The internal structure disappears; only effective behavior
  remains. This resembles renormalization in statistical physics.
\end{definition}

% ── C.15 Fundamental Theorem ──────────────────────────────────
\section*{C.15 The Spherepop Fundamental Theorem}

\begin{theorem}{Spherepop Fundamental Theorem}{sp-fundamental}
  Every finite Spherepop expression defines a terminating constraint-descent
  flow whose terminal state is an equivalence-class representative in the
  quotient space of semantic evaluation.
\end{theorem}

\begin{proof}
  Termination follows from \cref{thm:sp-termination}.
  Constraint descent follows from Proposition~\ref{prop:eval-decreases}.
  The quotient construction: define $E_1 \sim E_2$ iff
  $\mathrm{eval}(E_1) = \mathrm{eval}(E_2)$; the terminal value is
  a representative of $[E]$ in the quotient.
\end{proof}

\medskip
\noindent
\textit{Appendix~D} will show that Boolean logic can be represented as a
Spherepop system whose nested regions correspond to logical operators.
